\fancyhead[L]{Methodik} % 1000 Wörter (reste 3995 mots)



\section{Methodik}
Um die Fragestellung dieser Arbeit zu beantworten, wurden alle Vorschläge zu den Bürgerhaushalten in Toulouse und Stuttgart aus dem Zeitraum 2021 bis 2023 erhoben. 
Die Analyse der in den Beiträgen erwähnten Themen erfolgt durch den Einsatz von \textit{structural topic modelling}.
Diese innovative Methode des \textit{natural language processing} ermöglicht die \textit{unsupervised} und theorielose Extraktion strukturierender Themen aus umfangreichen Textcorpora.

\subsection{Fallauswahl}

\subsection{Daten}

\subsubsection{Datengrundlage - Bürgerhaushalten in Toulouse und Stuttgart}

Eine wichtige 

=> Toulouse, site officiel "Je Participe"
    - Base de donnée pré-existante, mais incomplÈte

=> Stuttgart, site officiel Buergerhaushalt Stuttgart
    - Pas de base accessible, web-scraping de toutes les propositions par année avec BeautifulSoup sur Python. 



\subsection{\textit{Structural Topic Modelling}: Eine Methode der \textit{Natural Language Processing}}

Ende der 90er Jahre bestand im Bereich Information Retrieval (IR) ein Bedürfnis für einen Wandel, da Textcorpora zunehmend zu umfangreich für klassische qualitative Auswertungsmethoden wurden.
So entstanden ab Anfang der 2000er Jahre neue Modellierungsmethoden, die die Dimensionen längerer Textdokumente reduzieren können, ohne auf die Kerninformationen zu verzichten \parencite{Heiberger2025STMI}.
So entstand die Methode, die für diese Arbeit angewendet wird: \textit{Structural Topic Modelling}.


\subsubsection{\textit{Structural Topic Modelling} in der Sozial- und Politikforschung}

Anfang der 2000er Jahre wurden die ersten Topic Models entwickelt, die auch als \textit{Latent Dirichlet Allocation} (LDA) bezeichnet werden.
Das LDA-Modell ist ein dreistufiges bayesisches Modell. 
Es identifiziert „Bags-of-Words“, die über mehrere Dokumente hinweg gestreut sind, und bildet damit Themen (\glqq Topics\grqq{}). 
Die \textit{topic probabilities} eines jeden Dokuments sind eine explizite, vereinfachte Darstellung des Dokuments, die dessen Kerninformationen enthält (\cite[993]{Blei2003LatentDirichletAllocation} ; \cite[1171]{Heiberger2021FacetsSpecialization}).
LDA sind sogenannte \textit{unsupervised topic modelling}-Methoden, die auf eine bestimmte Wahrnehmung der Sprache hinweisen, als durch latente Dimensionen strukturiert, welche für die Akteure selbst nicht unbedingt bewusst sind.
So sind die Methoden des \textit{unsupervised topic modelling} besonders für explorative Verfahren und die Entdeckung von Sprachbereichen geeignet.
Jedoch existieren andere \textit{topic modelling} Methoden, welche teilweise oder völlig überwacht werden und andere Ziele erreichen können \parencite[611]{McFarland2013DifferentiatingLanguageUsage}.\\

Die vorliegende Arbeit stützt sich auf eine spezifische Art des \textit{topic modelling}, die von Roberts et al. \parencite*{Roberts2014StructuralTopicModels} entwickelt wurde: \textit{Structural Topic Modelling} (STM).
Diese Variante wurde speziell für die quantitative Analyse von Umfragefeldern mit offenen Antwortmöglichkeiten entworfen. 
Wie die anderen LDA-Modelle ermöglicht STM einen explorativen Ansatz, bei dem Themen nicht nur vermutet, sondern unabhängig von den theoretischen Erwartungen des Forschers entdeckt werden.
Allerdings ergänzt STM die LDA-Modelle, indem auch Dokumentvariabeln (\textit{docvars}) als Kovariaten im Modell berücksichtigt werden können, was Schätzung der \textit{Topics} verbessert.
Mit STM können somit nicht nur Themen entdeckt, sondern auch deren Wahrscheinlichkeiten in Bezug auf die Dokumentvariablen modelliert werden (\cite[1066]{Roberts2014StructuralTopicModels} ; \cite[1171]{Heiberger2021FacetsSpecialization}). 

Diese Modelle werden seit ihrer Entwicklung im Jahr 2014 regelmäßig im Forschungskontext verwendet. 
Auf der offiziellen Website der Entwickler dieser Methode \url{www.structuraltopicmodel.com}{www.structuraltopicmodel.com} werden über fünfzig verschiedene Forschungsbeiträge in der Politik-, Sozial- und Medienforschung aufgeführt, die \textit{structural topic modelling} als Methode verwenden.


\subsubsection{Umsetzung der STM Methode}

Um die erhobenen Textcorpora einander gegenüberstellen zu können, werden möglichst viele Ortsnamen systemtisch neutralisiert. 
Der Grund dafür ist, dass ähnliche Ortsnamen auf sehr verschiedene Realitäten hinweisen können: So gibt es in Toulouse eine \glqq{}\textit{rue Franz Schubert}\grqq{} und in Stuttgart eine \glqq Schubertstraße\grqq{}, die beide zu \glqq \textit{(Franz) Schubert Street übersetzt}\grqq{} werden, obwohl sie sich nicht besonders ähneln. 
Die STM würde in diesem Fall die beiden Straßen als identisch einstufen, was die Ergebnisse verfälschen kann, da der Straßenname kaum etwas über die Ähnlichkeit der beiden Straßen aussagt.
Für beide Städte wurden nicht nur Straßen im weiteren Sinne neutralisiert, sondern auch Schulen, Plätze, Wälder, Flüsse, Kirchen, Brücken, Friedhöfe und Parks.



In einem zweiten Schritt wurden die beiden Corpora mit dem R-Paket \glqq polyglotR\grqq{} alle Textabschnitte ins Englische übersetzt und in einem einzigen Corpus zusammengefügt.
Die Extraktion aller Tokens stützt sich auf das R-Paket \glqq quanteda\grqq{}. 
Für die Tokenisierung wurden nicht nur die einzelnen Wörter, sondern auch die 200 häufigsten Kombinationen aus zwei bis drei Wörtern (\textit{n-Grams}) berücksichtigt.
Aus diesen Tokens wurde dann eine Document-Term-Matrix (DTM) erstellt. 
Die DTM fasst für jeden Token zusammen, wie oft und in welchen Dokumenten des Corpus er verwendet wurde. 
Alle Tokens, die in weniger als zehn Angeboten vorkamen, wurden herausgefiltert und die Dokumentvariablen wurden gesetzt.
In diesem Fall sind die relevanten Dokumentvariablen die Stadt und das Jahr, in dem der Beitrag erstellt wurde, sowie die Anzahl der Zustimmungen, die der Beitrag erhalten hat.
Letztendlich werden die DTM und die Dokumentvariablen verwendet, um eine STM mithilfe des R-Pakets \glqq stm\grqq{}, welches von Roberts et al. 2019 entwickelt wurde, zu bauen.
Für diese Arbeit wurde der k-Wert auf 69 gesetzt, um Residuen zu minimieren und Exklusivität, \textit{lower-bound variation} und semantische Kohärenz zu maximieren.


=> stm-Paket

=> nachdem ...